\section{核心风险}
\begin{enumerate}
\item \textbf{周期反转持续性：}光纤量价齐升可能是供需阶段性改善，价格和毛利率若回落，H2利润假设会迅速下修。
\item \textbf{现金转换：}Q1经营现金流仍为负，应收账款和存货规模较高，利润兑现不等于现金兑现。
\item \textbf{财务资产波动：}云创数据公允价值变动已对H1税前利润造成约CNY0.91亿元拖累，归母利润可能继续偏离扣非利润。
\item \textbf{客户和订单缺口：}客户、订单金额、ASP、利用率、认证和交付节奏没有完整公开，不能用行业需求替代公司收入。
\item \textbf{负债与质押：}短期借款高于现金，控股股东存在较高质押，股价波动可能放大融资与控制权风险。
\item \textbf{竞争：}长飞、亨通、中天、烽火等综合厂商在规模、认证和运营商渠道上具有竞争，价格战会压缩毛利。
\item \textbf{资金拥挤：}7月14日成交额24亿元，融资资金和龙虎榜机构交替进出，事件预期落空时回撤速度可能快于基本面修复。
\end{enumerate}

\section{催化剂与失效条件}
\begin{exhibitbox}[表7-1\quad 监控清单]
\begin{tabularx}{\exhibitboxwidth}{L{2.8cm}X L{3.0cm}X}
\textbf{方向} & \textbf{正向催化剂} & \textbf{失效信号} & \textbf{估值动作} \\
H1正式半年报 & 扣非接近CNY2.80亿元中枢或以上 & 扣非低于CNY2.30亿元 & 下调基准分母 \\
光纤业务 & 分部收入、毛利、库存和订单同时改善 & 只有行业涨价，没有公司收入/毛利 & 不给予成长倍数 \\
现金流 & OCF随利润转正或明显收窄 & OCF持续恶化、应收账款加速 & 降低PE和情绪权重 \\
价格行为 & 回撤后缩量、基本面继续确认 & 跌破CNY18.50且利润预期下修 & 转为观察/减风险 \\
券商证据 & 原始研报披露目标价和完整预测 & 继续只有媒体转述 & Street权重保持0\% \\
\end{tabularx}
\sourcenote{analysis/valuation\_model.md；analysis/secondary\_market\_analysis.md；source\_exhaustion\_log.md}
\end{exhibitbox}

\section{投资建议}
对于已有仓位，建议把持仓逻辑从“AI/光纤题材重估”切换为“公司层盈利和现金流验证”，不在H1预告后的天量交易日追高。对于新增仓位，当前20.86元不具备本报告基准风险收益比，优先等待回撤至情绪锚附近或等待正式半年报把扣非利润、现金流、毛利率和光纤分部数据同时验证。若H1上沿兑现且H2现金流和分部利润继续改善，可向牛情景重新评估；若利润、现金或价格三者任一先破坏，则按观察/减风险处理。
